% !PRINTED 336
\weberpage{336}

Die Co\"efficienten der cubischen Gleichung sind bekannt, sobald $m^3, mr^2, r^6$, oder auch, wenn $m$ und $r^2$ bekannt sind.

Ist dann die Gleichung (14) gelöst, so sind $y_1,y_2,y_3$ bis auf dritte Einheitswurzeln, die aber der Relation $y_1y_2y_3=R$ genügen müssen, bestimmt, und weiter können auch diese Grössen der Natur der Sache nach nicht bestimmt werden.

Wenn die Curve $f$ eine reelle Curve ist, so ergiebt sich aus den Formeln (13) und (15) noch der Satz, dass, \gesperrt{wenn $m$ und $r$ reell sind, $D_1$ positiv ist, und folglich die Wurzeln der Gleichung (14) alle drei reell sind.}

Drückt man die Discriminante $D_1$ durch die Coëfficienten der Gleichung $P,Q,R$ aus, so fliesst aus der Formel (15) noch das Resultat, dass man das Quadrat der Covariante $K$ als ganze rationale Function der Covarianten $f,\Delta,J$ ausdrücken kann. Der Ausdruck wird ziemlich complicirt und soll hier nicht weiter verfolgt werden.

Um aber die biquadratische Gleichung zu bilden, von der $m$ und $r$ abhängen, ist es n\"othig, auf die Invarianten dritter Ordnung näher einzugehen.

\begin{center}
\S.\,91.

\gesperrt{Die Invarianten der Curve dritter Ordnung und die biquadratische Gleichung.}
\end{center}

Alle Curven dritter Ordnung, deren Gleichung in der canonischen Form durch dieselben Grössen $\Sigma y_1^3$, $y_1y_2y_3$ linear dargestellt werden kann, haben dieselben Punkte zu Wendepunkten, beispielsweise also die Curven $f=0$ und $\Delta=0$, und allgemeiner alle Curven der Schaar
\weq{1}{F = \lambda f + 6\mu\Delta = 0,}
wo $\lambda,\mu$ beliebige Parameter sind. Alle Curven dieser Schaar schneiden sich in neun festen Punkten, und eine solche Schaar wird ein \gesperrt{Curvenbüschel} genannt. Die neun festen Punkte, die für alle Curven des Büschels die Wendepunkte sind, heissen die \gesperrt{Grundpunkte}. Bilden wir von der Form $F$ die \gesperrt{Hesse'sche} Determinante, so wird auch diese Form linear durch $\Sigma y_1^3$ und $y_1y_2y_3$ ausgedrückt und kann daher nach \S.\,90, (7) und (8) auch linear durch $f$ und $\Delta$ ausgedrückt werden, d. h. wir erhalten eine Form desselben Büschels.
